% !TeX program = xelatex
\documentclass[12pt]{article}

\usepackage[a4paper,margin=2.5cm]{geometry}
\usepackage{amsmath,amssymb,bm}
\usepackage{booktabs,tabularx,array}
\usepackage{graphicx}
\usepackage{fontspec}
\usepackage[unicode,hidelinks]{hyperref}

\defaultfontfeatures{Ligatures=TeX}
\setmainfont[AutoFakeSlant=0.2]{Noto Serif CJK SC}
\setsansfont{Noto Sans CJK SC}
\setmonofont{Noto Sans Mono CJK SC}
\XeTeXlinebreaklocale "zh"
\XeTeXlinebreakskip=0pt plus 1pt minus 0.1pt
\setlength{\parindent}{2em}
\setlength{\parskip}{0.25em}
\setlength{\emergencystretch}{3em}
\allowdisplaybreaks
\graphicspath{{manuscript/figures/}{figures/}}

\newcommand{\ind}{\mathbb I}
\newcommand{\softmax}{\operatorname{softmax}}
\newcommand{\logit}{\operatorname{logit}}
\newcommand{\sigmoid}{\operatorname{sigmoid}}

\title{生成式有限规则搜索模型}
\author{}
\date{}

\begin{document}
\maketitle

\section{模型概述}

模型的核心假设是，学习者不会在每个试次同时比较所有可能规则，也不会始终沿用一条固定规则。相反，学习者只在容量有限的工作空间中保留少量候选规则，并根据反馈不断评估和替换这些规则。每条活跃规则都有两个随试次变化的量：工作空间规则信念表示在当前候选中相对更相信哪条规则，动态规则精度表示该规则给出的类别预测有多明确。

在当前试次，模型先把物理刺激转化为带有被试特异噪声的内部知觉。每条活跃规则再根据内部刺激与其类别边界的关系给出类别预测。模型允许两种方式把这些规则预测转化为最终选择。若不设置持续执行规则，所有活跃规则按照作答前信念共同决定选择；若设置持续执行规则，工作空间仍保留多个候选，但只有一条被保护的执行规则直接控制当前作答。两种执行结构共用相同的规则空间、搜索控制、记忆更新和动态规则精度，但在可替换槽位与最终选择读出上有所不同。

作答并收到反馈后，模型更新活跃规则的信念、有限记忆和动态规则精度。反馈还通过两个步骤控制下一试次的搜索。第一步决定是否搜索：最近一次试次的正确或错误直接设定基线搜索概率，更早积累的失败则在此基础上进一步提高搜索概率。第二步调节搜索范围：如果搜索已经发生，包含最近一次反馈的累计失败量决定搜索更偏向局部还是全局；累计失败较低时更倾向于在当前规则附近寻找相似规则，累计失败较高时更倾向于从整个规则目录寻找新规则。发生搜索时，模型淘汰部分旧规则、引入新规则，并把原有信念迁移到更新后的工作空间。所有这些更新都从下一试次开始影响选择。

知觉噪声、当前工作空间、规则替换和执行切换都无法从行为中直接观察。粒子滤波（particle filtering, PF）因此并行追踪多条可能的隐藏路径，每条路径代表一种可能的内部学习过程。观察到实际选择后，PF 提高与该选择更相容的路径权重，并降低其他路径的权重，由此得到边际选择概率和在线潜在状态估计。PF 是研究者用来推断隐藏过程的数值方法，不是额外的被试认知机制。

下文先以 condition 1 的二分类、二值反馈说明基础模型。Condition 2 已有四分类二值反馈的程序扩展，其规则数和类别概率维度应使用对应配置。Condition 3 除四分类外，还包括部分正确反馈和事先告知的层级知识；第 \ref{sec:condition3-adaptation} 节完整说明这部分待实现的适配方案。该节中的配对状态和更新公式属于拟议扩展，不表示现有程序已经完成迁移或拟合验证。

\section{任务输入、候选规则与潜在状态}

\subsection{试次层观测量}

模型在每个试次直接观察三个量。对被试 $s$ 的第 $t$ 个试次，物理刺激、类别选择和正确性反馈分别记为
\begin{equation}
\bm x_{s,t}\in[0,1]^4,
\qquad
y_{s,t}\in\{1,2\},
\qquad
r_{s,t}\in\{0,1\}.
\label{eq:observations}
\end{equation}
其中，$\bm x_{s,t}$ 是归一化后的四维刺激，$y_{s,t}$ 是被试选择的类别，$r_{s,t}=1$ 和 $r_{s,t}=0$ 分别表示选择正确和错误。下文省略被试下标 $s$。拟合时，刺激用于计算当前选择概率，实际选择用于修正粒子权重，选择与反馈随后共同更新潜在状态。口头报告和反应时不进入模型，只在拟合后作为外部验证变量。

\subsection{固定候选规则目录}

模型假设，学习者可能考虑的规则来自一个固定目录。该目录包含 29 条规则，记为
\begin{equation}
\mathcal H=\{1,\ldots,J\},
\qquad
J=4+6+6+3+6+4=29.
\label{eq:rule-space}
\end{equation}
每条规则都把四维刺激空间划分为两个类别区域。目录包括六种几何形式：4 条单维阈值规则 $x_i=0.5$，6 条成对顺序规则 $x_i=x_j$，6 条成对和阈值规则 $x_i+x_j=1$，3 条两组和顺序规则 $x_i+x_j=x_k+x_l$，6 条成对近似规则，以及 4 条单维中等值规则。后两类规则用容差为 $0.10$ 的带状区域表示“两个维度足够接近”或“某个维度接近中间值”：
\begin{align}
R_{h,1}^{\mathrm{pair}}
&=\{\bm x:|x_i-x_j|\leq0.10\},
&
R_{h,2}^{\mathrm{pair}}
&=[0,1]^4\setminus R_{h,1}^{\mathrm{pair}},
\nonumber\\
R_{h,1}^{\mathrm{center}}
&=\{\bm x:|x_i-0.5|\leq0.10\},
&
R_{h,2}^{\mathrm{center}}
&=[0,1]^4\setminus R_{h,1}^{\mathrm{center}}.
\label{eq:band-rules}
\end{align}
每种几何划分还需要规定边界哪一侧对应类别 1。若把两个相反方向都作为独立规则，模型便难以区分几何规则和类别方向。为减少这种不可辨识性，主模型为每个几何划分固定一个与任务一致的类别方向，不再单独模拟方向学习。这是经过独立敏感性分析后采用的建模简化，并不表示任务已经把类别方向告诉被试。

在观察被试行为之前，模型不偏好任何规则，因此基础先验为
\begin{equation}
p_0(h)=\frac{1}{29}.
\label{eq:base-prior}
\end{equation}

\subsection{有限工作空间与潜在状态}

29 条规则构成完整搜索目录，但被试 $s$ 在任一试次只保留 $M_s$ 条活跃规则。$M_s$ 是离散的被试层容量参数，并在该被试的全部试次中保持不变。第 $t$ 个试次的活跃规则集合 $A_t$ 满足
\begin{equation}
A_t\subset\mathcal H,
\qquad
|A_t|=M_s.
\label{eq:active-set}
\end{equation}
$M_s$ 表示始终填满的固定槽位数，而不是规则数上限。因此，搜索只替换规则，不改变工作空间大小。第一试次按照 $p_0$ 不放回抽取 $M_s$ 条规则，并在抽中集合内重新归一化信念。基准配置取 $M_s=3$；正式拟合可从预先规定的离散候选中为不同被试选择容量。

二元结构参数 $\chi_s\in\{0,1\}$ 表示被试是否使用持续执行规则。当 $\chi_s=1$ 时，粒子还保存一条当前执行规则
\begin{equation}
e_t\in A_t,
\label{eq:executed-rule}
\end{equation}
其中 $e_t$ 直接控制当前作答。当 $\chi_s=0$ 时，不定义额外执行规则，所有活跃规则按照作答前信念共同参与选择。$\chi_s$ 在同一被试内保持不变，但不同被试可以采用不同结构。

给定被试参数 $\theta_s$，每个粒子保存一种可能的当前内部状态：
\begin{equation}
z_t=
\bigl(
\widetilde{\bm x}_t,
A_t,
\pi_t^{-},
F_t,
\bm\beta_t,
e_t\ind(\chi_s=1)
\bigr),
\label{eq:latent-state}
\end{equation}
其中，$\widetilde{\bm x}_t$ 是内部感知刺激，$A_t$ 是活跃规则集合，$\pi_t^{-}$ 是作答前的工作空间规则信念。$F_t$ 记录进入当前试次时可用的失败累积量，$\bm\beta_t$ 保存各条活跃规则的动态规则精度；仅当 $\chi_s=1$ 时，状态中才包含执行规则 $e_t$。这些变量共同描述被试看到了什么、正在考虑和相信哪些规则、此前积累了多少失败，以及当前怎样产生选择。


\section{从刺激到单条规则的类别预测}

\subsection{从物理刺激到内部知觉}

模型不假定被试能够无误差地感知物理刺激。对刺激的第 $j$ 个维度，内部感知值为
\begin{equation}
\widetilde x_{t,j}
=\operatorname{clip}(x_{t,j}+\varepsilon_{s,t,j},0,1).
\label{eq:perception}
\end{equation}
其中，$\varepsilon_{s,t,j}$ 表示被试特异的知觉误差，$\operatorname{clip}$ 将结果限制在 $[0,1]$ 内。噪声参数来自正式学习任务之前的知觉测量，并按照学习任务中的实际特征顺序排列。若知觉测量给出误差均值和标准差，模型逐维使用正态噪声；若测量给出阈限范围，则使用以 0 为中心、半宽由阈限决定的均匀噪声。不同粒子使用彼此独立但可复现的噪声流，因此 PF 会对无法直接观察的内部刺激进行边缘化。

\subsection{单条规则的类别概率}

规则 $h$ 将刺激空间划分为两个类别区域 $R_{h,1}$ 和 $R_{h,2}$。模型先计算内部刺激到每个类别区域的最短欧氏距离：
\begin{equation}
d_{t,h,c}
=\inf_{\bm z\in R_{h,c}}
\lVert\widetilde{\bm x}_t-\bm z\rVert_2.
\label{eq:boundary-distance}
\end{equation}
如果内部刺激已经位于类别 $c$ 的区域内，则 $d_{t,h,c}=0$。如果一个类别区域由多个互不相连的部分组成，则取距离最近的部分。这样，模型只在规则原有边界附近引入概率性，而不会改变规则本身如何划分类别。

随后，动态规则精度 $\beta_t(h)$ 把两个距离转换为规则 $h$ 的类别预测概率：
\begin{equation}
p_{t,h}(c)
=
\frac{\exp[-\beta_t(h)d_{t,h,c}]}
{\sum_{c'=1}^{2}\exp[-\beta_t(h)d_{t,h,c'}]}.
\label{eq:rule-emission}
\end{equation}
$\beta_t(h)$ 越高，距离差异对预测的影响越强，规则给出的类别概率也越明确；$\beta_t(h)$ 越低，两个类别的概率越接近 $0.5$。模型只使用这一套动态规则精度。同一个 $\beta_t(h)$ 既决定按照规则作答时有多稳定，也决定反馈对该规则的诊断性有多强。因此，$\beta_t(h)$ 表示规则预测映射的精度，而不只是外显动作的确信度；模型不再设置第二个独立的证据精度或额外的选择锐化参数。

\subsection{选择与反馈提供的规则证据}

观察到选择 $y_t$ 和反馈 $r_t$ 后，模型检验每条活跃规则能否解释这一结果。规则 $h$ 获得的未归一化支持为
\begin{equation}
\ell_t(h)
=r_t p_{t,h}(y_t)
+(1-r_t)\bigl[1-p_{t,h}(y_t)\bigr].
\label{eq:feedback-support}
\end{equation}
如果选择正确，能够预测已选类别的规则获得较高支持；如果选择错误，倾向于预测另一个类别的规则获得较高支持。模型再把这些支持值在当前活跃规则内归一化：
\begin{equation}
L_t(h)
=
\frac{\ind(h\in A_t)\ell_t(h)}
{\sum_{u\in A_t}\ell_t(u)}.
\label{eq:normalized-likelihood}
\end{equation}
归一化后的 $L_t(h)$ 表示本次选择和反馈对规则 $h$ 的相对证据。该操作不改变活跃规则之间的似然比，只保证后续更新具有稳定的数值尺度。


\section{反馈证据驱动的规则信念更新}

规则信念只在当前工作空间内比较。作答前信念 $\pi_t^{-}$ 表示观察当前选择之前相对相信哪些规则，$L_t$ 表示本次选择和反馈提供的新证据。模型将两者结合，得到反馈后的规则信念：
\begin{equation}
\pi_t^{+}(h)
=
\frac{
\ind(h\in A_t)
[\pi_t^{-}(h)]^{\gamma_s}
[L_t(h)]^{\alpha}
}{
\sum_{u\in A_t}
[\pi_t^{-}(u)]^{\gamma_s}
[L_t(u)]^{\alpha}
},
\qquad
\alpha=1.
\label{eq:fading-memory}
\end{equation}
等价的对数形式为
\begin{equation}
\log\pi_t^{+}(h)
=\gamma_s\log\pi_t^{-}(h)+\alpha\log L_t(h)-\log Z_t.
\label{eq:fading-memory-log}
\end{equation}
$\gamma_s\in[0,1]$ 控制旧信念保留多少。$\gamma_s$ 越小，旧规则之间的信念差异被压缩得越快；$\gamma_s$ 越大，过去证据保留得越久；当 $\gamma_s=1$ 时，模型退化为没有额外遗忘的顺序贝叶斯更新。反馈增益固定为 $\alpha=1$，静态记忆权重固定为 0，因此模型只保留这一条会衰减的证据记忆，不再用第二条永久记忆通道解释规则信念变化。

如果搜索改变了工作空间，模型先按照式 \eqref{eq:similarity-transport} 把旧信念迁移到新工作空间，得到 $\pi_t^{-}$；观察选择和反馈后，再按照式 \eqref{eq:fading-memory} 更新为 $\pi_t^{+}$。因此，新规则的初始信念来自规则之间的结构关系，而不是来自任意设置的独立记忆值。


\section{反馈驱动的下一试次搜索}

\subsection{失败信息的跨试次累积}

搜索控制器只读取已经完成的试次。进入试次 $t$ 时，可用的失败累积量为
\begin{equation}
F_t
=\delta_F F_{t-1}+(1-\delta_F)(1-r_{t-1}),
\qquad
F_1=0,
\qquad
\delta_F=0.60.
\label{eq:failure-accumulator}
\end{equation}
上一试次错误时，$(1-r_{t-1})=1$，因此新的失败被加入 $F_t$；上一试次正确时，不再加入新失败，原有失败只按 $\delta_F=0.60$ 保留。连续错误会使 $F_t$ 逐步升高，随后连续正确则使其逐步衰减。模型只保存这一条反馈历史，不再额外维护“掌握程度”累积量。

\subsection{新规则搜索的发生概率}

模型先根据最近一次反馈设定基线搜索概率：
\begin{equation}
E_t^{\mathrm R}
=
\begin{cases}
E_{C,s}, & r_{t-1}=1,\\
E_{E,s}, & r_{t-1}=0,
\end{cases}
\qquad
E_{E,s}\geq E_{C,s},
\label{eq:reactive-event}
\end{equation}
其中，$E_{C,s}$ 和 $E_{E,s}$ 分别表示上一试次正确或错误后，至少替换一个工作空间槽位的基线概率，并要求错误后的搜索概率不低于正确之后。第一试次令 $E_1^{\mathrm R}=E_{C,s}$，从而不再增加单独的初始搜索参数。

更早的失败历史可以在这一基线上进一步提高搜索概率：
\begin{equation}
E_t
=\sigmoid\!\left[
\logit(E_t^{\mathrm R})+c_{A,s}F_{t-1}
\right],
\qquad
c_{A,s}\geq0,
\qquad t\geq2.
\label{eq:accumulated-event}
\end{equation}
第一试次固定使用 $E_1=E_{C,s}$，累计项为 0。对第 $t$ 个试次，最近一次反馈 $r_{t-1}$ 只通过 $E_t^{\mathrm R}$ 进入模型；$F_{t-1}$ 只包含截至试次 $t-2$ 的更早失败。因此，模型不会把最近一次错误重复计算两遍。

这个结构形成三个清楚的嵌套模型。当 $E_{C,s}=E_{E,s}$ 且 $c_{A,s}=0$ 时，搜索概率固定；当 $E_{C,s}\neq E_{E,s}$ 且 $c_{A,s}=0$ 时，搜索只响应最近一次反馈；当 $c_{A,s}>0$ 时，更早的失败历史才会在即时反馈之外提供额外影响。该滞后读出不增加新的状态或参数。

\subsection{局部与全局搜索范围}

如果搜索已经发生，模型再决定搜索更偏向当前规则附近，还是整个规则目录。全局搜索所占的比例为
\begin{equation}
g_t
=g_{0,s}+(1-g_{0,s})c_{G,s}F_t,
\qquad
0\leq g_{0,s},c_{G,s}\leq1.
\label{eq:global-search}
\end{equation}
$g_{0,s}$ 是没有累计失败时的全局搜索基线，$c_{G,s}$ 控制失败累积把搜索推向全局的程度。这里使用包含最近一次反馈的 $F_t$。当 $c_{G,s}=0$ 时，全局搜索比例固定为 $g_{0,s}$；当 $c_{G,s}>0$ 时，持续失败会逐渐扩大搜索范围。

因此，进入试次 $t$ 时有三种可能：不搜索、局部搜索或全局搜索。对应概率为
\begin{align}
P_t(\mathrm{exploit})&=1-E_t,
\nonumber\\
P_t(\mathrm{local\ explore})&=E_t(1-g_t),
\nonumber\\
P_t(\mathrm{global\ explore})&=E_tg_t.
\label{eq:strategy-probabilities}
\end{align}
这些概率由反馈历史逐试次生成，不是看到行为曲线后再人为划分的学习阶段。


\section{搜索后的规则替换}

\subsection{规则替换数量}

首先确定有多少槽位允许被替换。可搜索槽位数为
\begin{equation}
S_s=M_s-\chi_s.
\label{eq:searchable-slots}
\end{equation}
当 $\chi_s=0$ 时，全部 $M_s$ 个槽位都可以替换；当 $\chi_s=1$ 时，执行规则占用并保护一个槽位，只搜索其余 $M_s-1$ 个槽位。模型把整体搜索概率 $E_t$ 转换为单个槽位的替换率：
\begin{equation}
m_t=1-(1-E_t)^{1/S_s},
\qquad
K_t\sim\operatorname{Binomial}(S_s,m_t).
\label{eq:replacement-count}
\end{equation}
其中 $K_t$ 是本次实际替换的槽位数。上述变换保证 $P(K_t>0)=E_t$，因此“是否发生搜索”的含义不会因为工作空间容量或可搜索槽位数不同而改变。

接下来决定淘汰哪些规则。模型从可替换的活跃规则中不放回抽取 $K_t$ 条，抽取权重为
\begin{equation}
P(h\text{ 被移出}\mid A_{t-1})
\propto
1-\pi_{t-1}^{+}(h)+\epsilon,
\qquad
\epsilon=10^{-12}.
\label{eq:drop-probability}
\end{equation}
因此，反馈后信念越低的规则越容易被淘汰。$\epsilon$ 只用于避免所有抽样权重同时为 0；当 $\chi_s=1$ 时，受保护的执行规则不进入淘汰候选池。

\subsection{新规则的局部与全局来源}

局部搜索需要先定义“哪些规则彼此相似”。模型用两条固定标签规则在均匀刺激空间中给出相同类别的概率衡量功能相似度：
\begin{equation}
\operatorname{sim}(h,u)
=P_{\bm x\sim\operatorname{Unif}([0,1]^4)}
\bigl[c_h(\bm x)=c_u(\bm x)\bigr].
\label{eq:functional-similarity}
\end{equation}
实现使用 $100{,}000$ 个均匀样本预先计算并固定这一相似度矩阵。令 $d(h,u)=\operatorname{clip}[1-\operatorname{sim}(h,u),0,1]$，局部搜索核为
\begin{equation}
K_L(u\mid h)
\propto
p_0(u)\exp[-d(h,u)/\tau_L],
\qquad
u\neq h,
\qquad
\tau_L=0.10.
\label{eq:local-kernel}
\end{equation}
$\tau_L=0.10$ 在所有被试间固定，$g_t$ 因而成为唯一随被试和反馈变化的局部／全局混合控制量。这里的相似度只表示两条规则在任务中产生的类别输出有多一致，不等同于被试口头描述中的语义相似性。

设 $I_{t-1}=\mathcal H\setminus A_{t-1}$ 为当前不在工作空间中的规则。局部搜索、全局搜索以及二者混合后的新人提议分布分别为
\begin{align}
Q_{L,t}(u)
&\propto
\ind(u\in I_{t-1})
\sum_{h\in A_{t-1}}
\pi_{t-1}^{+}(h)K_L(u\mid h),
\nonumber\\
Q_{G,t}(u)
&\propto
\ind(u\in I_{t-1})p_0(u),
\nonumber\\
Q_t(u)
&=(1-g_t)Q_{L,t}(u)+g_tQ_{G,t}(u).
\label{eq:newcomer-proposal}
\end{align}
局部提议 $Q_{L,t}$ 以整个工作空间的反馈后信念为中心，信念越高的旧规则影响越大，而不是只围绕信念最高的一条规则搜索。全局提议 $Q_{G,t}$ 则回到 29 条规则的均匀基础先验。$g_t$ 将二者混合为 $Q_t$，模型最后从中不放回抽取 $K_t$ 条新规则。

\subsection{工作空间更新后的信念迁移}

引入新规则后，模型还要为更新后的工作空间分配作答前信念。它不会把某条被淘汰规则的信念任意复制给与之配对的新规则。令保留规则集合和新进入规则集合分别为
\begin{equation}
\mathcal S_t=A_{t-1}\cap A_t,
\qquad
\mathcal N_t=A_t\setminus A_{t-1},
\qquad
q_t=\frac{|\mathcal N_t|}{M_s}=\frac{K_t}{M_s}.
\label{eq:transport-fraction}
\end{equation}
其中 $q_t$ 是被替换槽位占全部工作空间的比例。旧信念在保留规则上的归一化分量为
\begin{equation}
\Gamma_t(u)
=
\frac{
\ind(u\in\mathcal S_t)\pi_{t-1}^{+}(u)
}{
\sum_{v\in\mathcal S_t}\pi_{t-1}^{+}(v)
}.
\label{eq:survivor-carryover}
\end{equation}
若没有任何规则被保留，则定义 $\Gamma_t(u)=0$。旧工作空间信念向新工作空间的局部投影、全局投影和混合投影为
\begin{align}
Z_{L,t}(u)
&\propto
\ind(u\in A_t)
\sum_{h\in A_{t-1}}
\pi_{t-1}^{+}(h)K_L(u\mid h),
\nonumber\\
Z_{G,t}(u)
&\propto
\ind(u\in A_t)p_0(u),
\nonumber\\
Z_t(u)
&=(1-g_t)Z_{L,t}(u)+g_tZ_{G,t}(u).
\label{eq:semantic-projection}
\end{align}
主模型将保留分量和相似度投影按照槽位替换比例混合，得到新的作答前信念：
\begin{equation}
\boxed{
\pi_t^{-}(u)
=(1-q_t)\Gamma_t(u)+q_tZ_t(u)
}.
\label{eq:similarity-transport}
\end{equation}
迁移强度完全由实际替换比例 $K_t/M_s$ 决定，不增加新的拟合参数。如果 $K_t=0$，工作空间没有变化，直接令 $\pi_t^{-}=\pi_{t-1}^{+}$；如果全部槽位都被替换，新的信念完全由相似度投影 $Z_t$ 决定。

式 \eqref{eq:similarity-transport} 不只是数值处理，还表达一个认知假设：替换多少工作空间槽位，就按相同比例重新分配规则信念。因此，即使被淘汰规则原本只有很低的信念，更换一个槽位仍可能明显改变整个工作空间的信念分布。

为检验结论是否依赖这种“槽位比例重分配”，模型审计加入一个不增加自由参数的质量守恒信念迁移反事实。令 $\mathcal D_t=A_{t-1}\setminus A_t$，被淘汰规则原有的总信念质量为
\begin{equation}
m_t^{\mathrm{drop}}
=\sum_{h\in\mathcal D_t}\pi_{t-1}^{+}(h),
\label{eq:dropped-belief-mass}
\end{equation}
再将式 \eqref{eq:semantic-projection} 限制在新进入规则集合 $\mathcal N_t$ 上并重新归一化，得到 $\widetilde Z_t$。质量守恒反事实定义为
\begin{equation}
\pi_{t,\mathrm{MP}}^{-}(u)
=
\begin{cases}
\pi_{t-1}^{+}(u),&u\in\mathcal S_t,\\
m_t^{\mathrm{drop}}\widetilde Z_t(u),&u\in\mathcal N_t.
\end{cases}
\label{eq:mass-preserving-transport}
\end{equation}
在这一版本中，保留规则维持原来的绝对信念质量，新规则总共只继承被淘汰规则留下的质量。式 \eqref{eq:similarity-transport} 仍是主模型；式 \eqref{eq:mass-preserving-transport} 只用于一次性的选择负对数似然（NLL）、后验预测检验（PPC）和潜在状态敏感性分析，不作为新的被试层结构参数。

\subsection{搜索与执行规则切换}

当 $\chi_s=1$ 时，搜索只表示候选工作空间发生变化，并不意味着被试一定立刻改变实际执行的规则。只有在 $K_t>0$ 且独立切换事件发生时，执行规则才会改变：
\begin{equation}
P(e_t\neq e_{t-1}\mid K_t>0)=s_{\mathrm{exec}},
\qquad
s_{\mathrm{exec}}=0.20.
\label{eq:execution-switch}
\end{equation}
由于 $P(K_t>0)=E_t$，执行规则的边际切换概率为 $E_t s_{\mathrm{exec}}$。一旦发生切换，新执行规则从 $A_t\setminus\{e_{t-1}\}$ 中按照迁移后的作答前信念抽取。$s_{\mathrm{exec}}=0.20$ 在所有被试间固定，因此持续执行结构只增加二元参数 $\chi_s$，不再增加一个连续的执行灵活性参数。


\section{反馈驱动的动态规则精度更新}

每条活跃规则都维护一个动态规则精度 $\beta_t(h)$。观察选择和反馈后，模型先计算该结果对规则 $h$ 的支持程度：
\begin{equation}
b_t(h)
=r_t p_{t,h}(y_t)
+(1-r_t)[1-p_{t,h}(y_t)],
\qquad
\kappa_t(h)=2[b_t(h)-0.5].
\label{eq:beta-evidence}
\end{equation}
$b_t(h)$ 越大，当前结果与规则 $h$ 越一致；$\kappa_t(h)$ 将其重新缩放到 $[-1,1]$。$\kappa_t(h)>0$ 表示规则得到支持，$\kappa_t(h)<0$ 表示规则受到反驳。所有活跃规则随后按照同一反馈更新：
\begin{equation}
\beta_{t+1}(h)=
\begin{cases}
\min\!\left\{
\beta_t(h)+\eta_{+,s}\kappa_t(h)
[\beta_{\max}-\beta_t(h)],
\beta_{\max}
\right\},
&\kappa_t(h)\geq0,\\[6pt]
\max\!\left\{
\beta_t(h)-\eta_{-,s}\beta_t(h)
\min[1,-\kappa_t(h)],
\beta_{\min}
\right\},
&\kappa_t(h)<0,
\end{cases}
\label{eq:beta-update}
\end{equation}
其中
\begin{equation}
\beta_{\min}=0.1,
\qquad
\beta_{\max}=25.
\label{eq:beta-bounds}
\end{equation}
$\beta_{0,s}$、$\eta_{+,s}$ 和 $\eta_{-,s}$ 是被试层参数。规则得到支持时，精度按照距离上限 $\beta_{\max}$ 的剩余空间增加，所以越接近上限，增长越慢；规则受到反驳时，精度按照当前数值的一定比例下降，但不会低于 $\beta_{\min}$。

模型更新 $A_t$ 中的全部活跃规则，而不只更新执行规则 $e_t$。因此，即使 $\chi_s=1$，没有直接控制本次作答的候选规则也会接受同一结果的反事实评价。新进入或重新进入工作空间的规则都从 $\beta_{0,s}$ 开始；规则离开工作空间后，不再保存可供将来恢复的精度状态，也不根据新人先验额外调整初值。


\section{从工作空间规则到最终选择}

\subsection{无持续执行规则的选择}

当 $\chi_s=0$ 时，所有活跃规则共同参与选择。最终类别概率是规则条件预测按照作答前信念的加权平均：
\begin{equation}
p_t^{\mathrm{choice}}(c)
=\sum_{h\in A_t}\pi_t^{-}(h)p_{t,h}(c).
\label{eq:choice-execution-off}
\end{equation}
这里直接使用 $\pi_t^{-}$。模型不会对规则信念再次取幂，不会只选择信念最高的规则，也不会在每个试次额外抽取一条临时执行规则。

\subsection{持续执行规则下的选择}

当 $\chi_s=1$ 时，只有当前执行规则 $e_t$ 直接产生选择：
\begin{equation}
p_t^{\mathrm{choice}}(c)=p_{t,e_t}(c).
\label{eq:choice-execution-on}
\end{equation}
其余活跃规则仍会影响哪些规则被淘汰、哪些新规则被提出、信念如何迁移以及以后是否切换执行规则，但不会再次混入当前试次的选择概率。

\subsection{执行结构与输出噪声}

两种执行结构都把选择锐化幂次固定为 1，并关闭策略确信度锐化、规则承诺锐化和输出 lapse（随机失误）。因此，模型输出的选择概率直接等于上述认知过程产生的概率：
\begin{equation}
p_t^{\mathrm{obs}}(c)=p_t^{\mathrm{choice}}(c).
\label{eq:no-output-noise}
\end{equation}
因此，预测随试次发生的变化只能来自知觉、规则信念、工作空间搜索、执行结构或动态规则精度，不能被一个独立的输出噪声机制吸收。


\section{单试次的生成顺序}

每个试次严格分为两个阶段：模型先生成作答前预测，观察到选择和反馈后再学习。
\begin{equation}
\operatorname{begin\_trial}(\bm x_t)
\longrightarrow
p_t^{\mathrm{obs}}(y)
\longrightarrow
\operatorname{complete\_trial}(y_t,r_t).
\label{eq:lifecycle}
\end{equation}
对每个粒子，完整顺序为
\begin{equation}
\begin{aligned}
&\pi_{t-1}^{+},F_{t-1},\bm\beta_t,A_{t-1},e_{t-1}
\\
&\quad\longrightarrow
\widetilde{\bm x}_t,F_t,E_t,g_t,K_t,A_t,\pi_t^{-},e_t
\\
&\quad\longrightarrow
p_t^{\mathrm{obs}}(y)
\\
&\quad\longrightarrow
\text{观察 }(y_t,r_t)
\longrightarrow
L_t,\pi_t^{+},\bm\beta_{t+1}.
\end{aligned}
\label{eq:causal-order}
\end{equation}
第一行是上一试次结束后保留下来的状态。第二行表示模型在看到当前选择之前生成内部知觉、搜索决策、新工作空间和作答前信念。第三行是已经可以对外输出的选择概率。只有在这之后，模型才观察 $y_t$ 和 $r_t$，并更新规则证据、反馈后信念和下一试次使用的动态规则精度。

因此，当前试次的选择和反馈不能反向改变本试次已经输出的预测。反馈 $r_t$ 对搜索控制器的影响从试次 $t+1$ 才开始。这一时间边界避免模型在预测当前选择时提前使用尚未发生的信息。


\section{粒子滤波对潜在路径的汇总}

一个粒子只代表一种可能的内部学习路径，而行为数据通常同时支持多种路径。PF 使用 $R$ 个粒子并行表示这些可能性。设试次 $t$ 作答前第 $i$ 个粒子的权重为 $w_{t-1}^{(i)}$，该粒子独立运行生成过程并得到 $p_t^{(i)}(c)$。PF 将所有粒子的预测按权重平均，得到作答前边际选择概率：
\begin{equation}
\widehat p_t(c)
=\sum_{i=1}^{R}w_{t-1}^{(i)}p_t^{(i)}(c).
\label{eq:pf-predictive}
\end{equation}
观察到实际选择 $y_t$ 后，与该选择更相容的粒子获得更高权重：
\begin{equation}
\widetilde w_t^{(i)}
\propto
w_{t-1}^{(i)}p_t^{(i)}(y_t),
\qquad
\sum_i\widetilde w_t^{(i)}=1.
\label{eq:pf-weight-update}
\end{equation}
权重更新完成后，每个粒子再使用选择和反馈执行式 \eqref{eq:normalized-likelihood}、\eqref{eq:fading-memory} 和 \eqref{eq:beta-update}。粒子权重的集中程度用有效样本量表示：
\begin{equation}
\operatorname{ESS}_t
=\frac{1}{\sum_i(\widetilde w_t^{(i)})^2}.
\label{eq:ess}
\end{equation}
当 $\operatorname{ESS}_t<R/2$ 时，PF 进行系统重采样：保留更多高权重路径，减少低权重路径，并为复制后的粒子重新分配独立的未来随机数流。重采样复制完整状态，包括工作空间、规则信念、记忆、动态规则精度和可选执行规则。重采样只是研究者维持数值精度的方法，不是被试发生了一次额外的认知更新。

PF 输出需要区分两类时间点。作答前预测量（predictive）只使用截至试次 $t-1$ 的行为历史；过滤后量（filtered）则在观察当前选择 $y_t$ 后重新加权。二者都是在线过滤结果，不会利用未来试次重新解释较早状态，因此都不属于事后平滑（smoothing）。在固定被试参数 $\theta_s$ 的条件下，PF 表示的是潜在路径不确定性，而不是参数 $\theta_s$ 的后验不确定性。

为减少 PF 随机性，同一参数点可以使用 $B$ 个独立随机种子。模型先在每个试次平均边际选择概率：
\begin{equation}
\overline p_t(c)=\frac{1}{B}\sum_{b=1}^{B}\widehat p_t^{(b)}(c),
\label{eq:seed-aggregation}
\end{equation}
然后再由 $\overline p_t(y_t)$ 计算选择负对数似然。模型平均的是逐试次选择概率，而不是先分别计算多个 NLL 再取平均。粒子数 $R$、随机种子数 $B$ 和重采样阈值都属于数值精度设置，不能解释为被试层认知参数。


\section{个体参数的估计}
\label{sec:individual-estimation}

前述生成过程回答的是：给定一组参数，被试可能怎样学习和选择。个体参数估计反过来问：在这些可能的参数中，哪一组最能解释这名被试实际做出的选择。这里对每名被试分别拟合一组在整个实验中保持不变的参数；随试次变化的是工作空间、规则信念和动态规则精度等内部状态。程序中的 Hyper-CD 指分块坐标搜索，并不表示这里采用了跨被试的层级贝叶斯估计。

\subsection{哪些量需要估计}

在完整 PMH 模型中，被试 $s$ 的待估计参数写为
\begin{equation}
\theta_s=(M_s,\chi_s,\gamma_s,E_{C,s},\delta_{E,s},
 g_{0,s},c_{A,s},c_{G,s},\beta_{0,s},\eta_{+,s},\eta_{-,s}).
\label{eq:subject-parameter-vector}
\end{equation}
其中，$M_s$ 与 $\chi_s$ 决定工作空间容量和执行结构，其余参数控制记忆、搜索与规则精度更新。感知参数由已有的个体知觉测量确定，在本轮选择拟合中固定；表 \ref{tab:parameter-scope} 中的共同设置也不随候选参数改变。因此，程序不会为了改善选择拟合而同时重新调整所有模型常数。

正确与错误后的搜索概率不是任意组合。搜索时使用 $E_{C,s}$ 和非负差值参数 $\delta_{E,s}$，再计算
\begin{equation}
E_{E,s}=\sigmoid\!\left[\logit(E_{C,s})+\delta_{E,s}\right].
\label{eq:fit-event-parameterization}
\end{equation}
这样自动满足 $E_{E,s}\geq E_{C,s}$：$\delta_{E,s}=0$ 表示两类反馈后的基础搜索概率相同，正值表示错误后更容易搜索。$c_{A,s}$ 和 $c_{G,s}$ 也保留精确的零候选，允许数据选择不需要相应失败累积增益的情况。零值具有明确的机制含义，不能用一个很小的正数代替。

\subsection{用实际选择的预测概率评价一组参数}

给定候选 $\theta_s$，程序从实验开始重新运行 PF，依次读入该被试的刺激、实际选择和反馈。试次 $t$ 的选择概率必须在读入当前选择与反馈之前形成；记为 $\overline p_t(c\mid\bm x_t,\mathcal D_{<t},\theta_s)$，其中 $\mathcal D_{<t}$ 是此前的行为历史。该概率按式 \eqref{eq:pf-predictive} 汇总粒子，再按式 \eqref{eq:seed-aggregation} 汇总独立 PF 重复。这里沿着被试真实经历的历史计算预测，不让模型自己生成的选择替换观察到的选择。

令 $\mathcal I_s$ 为参与参数评分的试次集合，$N_s=|\mathcal I_s|$。当前搜索使用平均选择负对数似然：
\begin{equation}
\mathcal J_s(\theta_s)
=-\frac{1}{N_s}\sum_{t\in\mathcal I_s}
\log\!\left[\max\!\left\{\overline p_t(y_t\mid\bm x_t,\mathcal D_{<t},\theta_s),10^{-12}\right\}\right].
\label{eq:individual-choice-objective}
\end{equation}
模型给实际选择的概率越高，损失越小。例如，把实际选择的概率从 $0.2$ 提高到 $0.8$，该试次的损失就从约 $1.61$ 降到 $0.22$。这要求模型解释被试究竟选了什么，包括错误选择，而不仅是追随一条平滑后的正确率曲线。$10^{-12}$ 只是避免数值上取 $\log 0$ 的下限，不是新增的输出噪声。

目前代表性被试的描述性拟合使用全部可用序列，公共指标的默认规则不计首试次，因此完整有效序列通常有 $N_s=T_s-1$。不计分的首试次仍参与状态更新。若进行时间前缀训练和后缀检验，则 $\mathcal I_s$ 只包含训练部分，最终复评分也沿用同一个集合，不能用后缀表现挑参数。同一被试、同一评分集合下，总 NLL 等于 $N_s\mathcal J_s$，二者给出的候选排序相同；停止阈值则必须与实际使用的平均 NLL 尺度一致。

当前目标只使用选择 NLL，没有加入反应时、口述一致性或参数先验惩罚。其估计可理解为在给定候选范围内，用数值近似的选择似然寻找较优参数，而不是直接抽样参数后验。尤其不能把 $B$ 次 PF 的 NLL 先平均来代替式 \eqref{eq:individual-choice-objective}：先平均概率与先取对数是两个不同的计算。

\subsection{把参数分成便于比较的八个坐标块}

全部参数组合逐一枚举代价很高，因此搜索每次只改变一个参数或一组相关参数，其他参数保持为当前取值。当前完整模型使用八个坐标块，按表 \ref{tab:fit-coordinate-blocks} 的顺序循环比较。

\begin{table}[htbp]
\centering
\small
\caption{完整模型的分块坐标及当前代表性拟合的候选数。粗、细候选数是每个坐标块的数量，不是全部联合组合数。}
\label{tab:fit-coordinate-blocks}
\begin{tabularx}{\textwidth}{l>{\raggedright\arraybackslash}Xrr}
\toprule
坐标块 & 一次共同改变的量 & 粗搜 & 细搜 \\
\midrule
工作空间与执行 & $(M_s,\chi_s)$ & 9 & 9 \\
记忆保持 & $\gamma_s$ & 7 & 13 \\
反应式搜索 & $(E_{C,s},\delta_{E,s})$，据此导出 $E_{E,s}$ & 30 & 99 \\
全局搜索 & $(g_{0,s},c_{G,s})$ & 42 & 143 \\
搜索事件的累积增益 & $c_{A,s}$ & 7 & 13 \\
初始规则精度 & $\beta_{0,s}$ & 6 & 11 \\
支持反馈更新率 & $\eta_{+,s}$ & 7 & 13 \\
反驳反馈更新率 & $\eta_{-,s}$ & 7 & 13 \\
\bottomrule
\end{tabularx}
\end{table}

例如，在比较记忆保持率时，容量、执行结构和其他参数都先固定，程序分别评价各个 $\gamma_s$ 候选；选定较优值之后，再比较反应式搜索参数。相关参数打包比较，是为了允许它们一起改变：容量与执行结构一起比较，正确/错误后的搜索概率通过式 \eqref{eq:fit-event-parameterization} 一起改变，全局搜索的基线与增益也一起比较。

当前容量候选为 $1$ 至 $5$。容量为 $1$ 时只有一条活跃规则，信念平均与持续执行的读出没有可区分的规则选择，因此只保留 $\chi_s=0$；其余容量各保留两种执行结构，共有九个组合。连续参数在模型定义上可以连续取值，但本轮算法只比较配置中预先给定的有限候选点，不在每个区间内做连续梯度优化。细搜采用预先加密的候选表，也不是围绕赢家无限缩小步长。

\subsection{先粗搜，再从多个较优位置细搜}

粗搜的任务是较快地找到值得继续比较的参数区域。当前代表性拟合从两个起点出发，分别开启和关闭持续执行，其余起始值相同。每个起点按八个坐标块依次比较一遍，称为一轮；改变前面坐标后，后面坐标会基于更新后的完整参数重新评价，因此下一轮仍可能进一步改善结果。

在单个坐标块内，程序评价该块的全部候选，并保留原参数点作为比较依据。只有平均 NLL 的下降达到 $10^{-4}$ 才接受移动，避免因极小变化不断切换参数。每个起点最多执行六轮；连续两轮没有接受任何改善就提前结束。这个停止条件只表示当前坐标搜索不再明显改善，不证明已经找到整个联合参数空间的全局最优解。

粗搜结束后，程序从已经评价过的不同完整参数点中取排名靠前的四个作为细搜起点。细搜使用更密的候选表和更高的 PF 预算，重新评价这些起点，再按同样的坐标顺序与停止规则搜索。四个起点可以来自同一次粗搜路径，不要求分别对应四种机制。保留多个起点的目的，是减少只从一个位置开始而遗漏其他较优区域的风险。

各阶段的数值预算不同：当前粗搜使用 $R=16,B=4$，细搜使用 $R=64,B=8$。它们改变的是概率估计的精细程度，不改变被试的试次数、认知机制或评分定义。低预算有可能把好参数排得较低，因此需要另行检查它能否保留高预算下的较优候选；增加搜索速度本身不构成数值可靠性的证明。

\subsection{用共同随机数比较，再用独立随机数定稿}

同一个参数点在不同 PF 随机种子下会有小幅评分变化。为使候选比较尽量反映参数差异，同一被试、同一搜索阶段的候选使用共同的 PF 随机种子；而同一参数点内部的 $B$ 次重复仍相互独立。共同种子并不要求不同参数生成完全相同的内部路径，也不能消除全部数值误差。

细搜完成后，程序对该阶段已评价的完整参数点去重，按评分选出前四个，组成最终候选集合 $\mathcal C_s$。这些候选使用与搜索阶段不同的一组随机种子，在 $R=128,B=16$ 的预算下重新运行。最终候选之间仍共享这组新种子，并仍然先平均逐试次概率，再计算式 \eqref{eq:individual-choice-objective}。最终估计为
\begin{equation}
\widehat\theta_s
=\underset{\theta\in\mathcal C_s}{\arg\min}\;
\mathcal J_s^{\mathrm{final}}(\theta).
\label{eq:final-subject-estimate}
\end{equation}
因此，低预算细搜的第一名不一定是最终采用的参数。独立复评分减少了某个候选仅因搜索种子有利而胜出的风险，但只能在保留下来的候选中选择，无法补回已经被搜索遗漏的区域。若候选评分相同，程序使用固定的排序规则保持结果可复现；这种顺序本身不表示数据更支持其中一个机制。

\subsection{并行、保存与拟合后的使用}

同一坐标块内，不同候选和不同 PF 重复可以分配到多个 CPU 进程；坐标之间则有先后依赖，必须等当前坐标比较结束，才能确定下一坐标要围绕哪组参数搜索。并行只改变计算分工，不改变候选的随机种子、评分集合和接受规则。已计算过的同阶段参数点由缓存复用，每完成一个坐标保存检查点；续跑前核对配置与数据等运行指纹，避免把不同任务的评分混在一起。

每名被试需要同时保存搜索阶段最优点、最终复评分最优点、各候选评分、各起点的坐标移动与停止原因，以及候选表、随机种子和 PF 预算。选定 $\widehat\theta_s$ 后，将参数固定并重新运行标准分析，输出选择预测和内部状态。后处理不能再依据口述对齐或图形是否好看改选参数，否则就改变了这里定义的估计目标。

这一流程给出的是一个选中的参数点及其候选比较记录。多个 PF 粒子表示该参数条件下的潜在路径不确定性，不是参数置信区间；多个搜索起点也不是被试参数的独立样本。即使状态递推没有使用未来反馈，只要参数由完整行为序列选出，同一序列上的拟合仍是样本内结果。参数是否可恢复、相近参数是否可区分，以及冻结参数后的留出预测是否可靠，需要另外检验，不能仅由搜索正常结束来判断。

上述候选数、停止设置和三阶段预算对应当前代表性被试的完整模型拟合配置，是明确的估计方案，不是认知模型不可更改的常数。复现某次估计时，以该次运行归档的配置为准；有关数值校准、恢复检验及论文呈现的安排另见配套文稿 \texttt{model\_0826\_plus.tex}。


\section{参数层级、数值设置与解释边界}

模型中的量分为三类：允许被试间不同的参数、所有被试共享的固定设置，以及只控制 PF 精度的数值设置。表 \ref{tab:parameter-scope} 汇总了这些层级。表中的基准值只是拟合前的配置起点，不代表所有被试的最终估计。

\begin{table}[htbp]
\centering
\footnotesize
\renewcommand{\arraystretch}{0.96}
\caption{模型参数的层级、作用和基准设置。}
\label{tab:parameter-scope}
\begin{tabularx}{\textwidth}{>{\raggedright\arraybackslash}p{0.19\textwidth}>{\raggedright\arraybackslash}p{0.21\textwidth}>{\raggedright\arraybackslash}X>{\raggedright\arraybackslash}p{0.17\textwidth}}
\toprule
参数 & 层级 & 作用 & 基准/约束 \\
\midrule
$M_s$ & 被试层离散 & 固定工作空间槽位数 & 基准为 3 \\
$\chi_s$ & 被试层二元 & 选择信念平均或持续执行规则 & $0$ 或 $1$ \\
$\gamma_s$ & 被试层连续 & 旧规则信念的保持程度 & 基准为 0.80 \\
$E_{C,s}$、$E_{E,s}$ & 被试层连续 & 正确/错误后的反应式搜索概率 & $0<E_C\leq E_E<1$ \\
$g_{0,s}$ & 被试层连续 & 全局新人提议的基线比例 & $[0,1]$ \\
$c_{A,s}$ & 被试层连续，含零边界 & 更早失败历史对搜索事件的额外增益 & $c_A\geq0$ \\
$c_{G,s}$ & 被试层连续，含零边界 & 失败累积对全局搜索的增益 & $[0,1]$ \\
$\beta_{0,s}$ & 被试层连续 & 初始及新人规则的动态规则精度 & 基准为 5 \\
$\eta_{+,s}$、$\eta_{-,s}$ & 被试层连续 & 支持/反驳反馈后的 $\beta$ 更新率 & 基准为 0.04/0.15 \\
$\delta_F$ & 共同固定 & 失败累积的保持率 & 0.60 \\
$\tau_L$ & 共同固定 & 局部规则相似度尺度 & 0.10 \\
$s_{\mathrm{exec}}$ & 共同固定 & 搜索发生后执行规则的切换比例 & 0.20 \\
$\beta_{\min},\beta_{\max}$ & 共同固定 & 动态规则精度的数值边界 & 0.1/25 \\
$\alpha,w_0$ & 共同固定 & 反馈增益与静态记忆权重 & 1/0 \\
$R,B$ & 数值设置 & PF 粒子数与独立随机种子数 & 由数值收敛确定 \\
\bottomrule
\end{tabularx}
\end{table}

模型边界同样由没有加入的机制决定。当前生成过程不包含类别标签方向学习、独立静态记忆通道、规则权重锐化、策略确信度锐化、输出 lapse、错误规则捕获或额外规则承诺机制。四分类类别排列、事后平滑、引导式粒子滤波以及替代的隐马尔可夫模型（HMM）或基于模拟的推断也不属于模型的一部分。

潜在状态的解释还受到工作空间和数值可靠性的限制。$\pi_t^{-}$ 和 $\pi_t^{+}$ 只是在当前粒子的有限工作空间内归一化的规则信念，不是 29 条规则上的无条件完整后验。跨粒子的过滤后（filtered）规则概率表示在给定模型和参数下，与截至当前试次的选择历史相容的在线状态分布。只有在粒子数和随机种子收敛、参数与状态恢复、生成数据检验以及独立行为指标验证都通过之后，才能把这些状态用于解释被试策略的形成、保持和切换。

\section{Condition 3 的层级反馈与按键配对学习}
\label{sec:condition3-adaptation}

Condition 3 的被试不仅要学习如何分类，还要逐渐理解哪些反应选项属于同一个科。模型因此需要同时学习分类规则和选项配对，使部分正确反馈的作用随已有知识而变化。下面说明这一扩展的主要机制；共享程序已实现这一版本，其科学有效性仍需通过拟合与恢复实验检验。

\subsection{被试知道什么，还需要学习什么}
\label{sec:condition3-coding}

任务的上层是科，下层是种。真实类别 1、2 属于一科，类别 3、4 属于另一科。Condition 3 的被试在实验前知道有两个科、每科两个种，但不知道哪两个按键属于同一科。模型应保留这份初始知识，同时让具体配对从学习过程中逐渐形成。

这里需要区分数据中的三个字段。\texttt{category} 是刺激的真实类别，\texttt{choice} 是被试所选答案的统一类别编号，\texttt{presskey} 是实际按键。每名被试的 \texttt{choice} 与 \texttt{presskey} 一一对应，但对应关系在被试之间不同。例如，301 的答案 1、2、3、4 分别对应按键 4、2、3、1，因此他的两个科实际上对应按键 $\{2,4\}$ 和 $\{1,3\}$。

这张对应表用于正确理解数据，不是要拟合的认知参数。模型可以继续用 \texttt{choice} 编号计算，但不能因为编号是 1 和 2 就预先认定它们同科。真实类别和正确配对由实验者掌握，用于生成反馈和评价结果；被试模型只能根据刺激、自己的反应和收到的反馈学习。

\subsection{分类规则与选项配对一起学习}
\label{sec:condition3-pair-state}

四个反应分成两组、每组两个，只有三种配对：
\begin{equation}
G_1=\{\{1,2\},\{3,4\}\},\quad
G_2=\{\{1,3\},\{2,4\}\},\quad
G_3=\{\{1,4\},\{2,3\}\}.
\label{eq:c3-pairings}
\end{equation}
这里使用统一的反应编号，展示个体结果时再换回实际按键。开始时，模型对三种配对同样相信；随后根据经验调整它们的权重。

配对判断会受到分类规则的影响。以 301 为例，按 4 得到 $0.5$ 时，如果某条规则认为这幅刺激应该按 2，那么反馈支持 2、4 同科；如果另一条规则认为应该按 3，同一反馈就可能支持 3、4 同科。因此，每条规则都需要保留自己的配对权重。

令 $\omega_t^{-}(a\mid h)$ 表示作答前，在规则 $h$ 下相信配对 $G_a$ 的程度。规则与配对的联合信念为
\begin{equation}
\Psi_t^{-}(h,a)=\pi_t^{-}(h)\omega_t^{-}(a\mid h),
\qquad
\sum_{a=1}^{3}\omega_t^{-}(a\mid h)=1.
\label{eq:c3-joint-prior}
\end{equation}
第一试次取 $\omega_1^{-}(a\mid h)=1/3$。这些权重随试次变化，不是额外拟合的三个常数参数。

第一版复用 condition 2 的 116 条四分类规则，基础先验在这些规则上均匀。单条规则仍按式 \eqref{eq:rule-emission} 计算类别概率，只将分母的求和范围改为四个反应。每条规则保留一个 $\beta$，工作空间容量仍为 $M_s$；增加的是对配对的判断，而不是三套独立的分类模型。

\subsection{三种反馈怎样更新规则和配对}
\label{sec:condition3-feedback}

收到反馈后，模型检查每种“规则与配对”的组合是否能解释这一结果。令 $s_a(y)$ 表示配对 $G_a$ 中与反应 $y$ 同科的另一个反应，则
\begin{equation}
\ell_t(h,a)=
\begin{cases}
p_{t,h}(y_t),&r_t=1,\\
p_{t,h}[s_a(y_t)],&r_t=0.5,\\
1-p_{t,h}(y_t)-p_{t,h}[s_a(y_t)],&r_t=0.
\end{cases}
\label{eq:c3-feedback-kernel}
\end{equation}
完全正确支持所选答案，部分正确支持同科的另一种，零分支持另一个科中的两个答案。三种情况互不重叠，概率之和为 1。因此，condition 3 的零分不能直接沿用二值反馈中的 $1-p_{t,h}(y_t)$，因为它还要排除同科另一种的概率。

记忆更新沿用前文“保留一部分旧信念，再加入新证据”的思路，将式 \eqref{eq:fading-memory} 扩展为
\begin{equation}
\Psi_t^{+}(h,a)=
\frac{[\Psi_t^{-}(h,a)]^{\gamma_s}\ell_t(h,a)}
{\sum_{u\in A_t}\sum_{b=1}^{3}[\Psi_t^{-}(u,b)]^{\gamma_s}\ell_t(u,b)}.
\label{eq:c3-joint-update}
\end{equation}
同一个 $\gamma_s$ 控制规则与配对证据的保留程度，不再设置配对专用的遗忘率。$\gamma_s=1$ 时完整保留旧证据，$\gamma_s=0$ 时把当前工作空间内的旧联合信念视为等权，再由本次反馈更新。这里衰减的是联合信念，与分别衰减规则权重和配对权重是不同的记忆假设。

更新后，把某条规则下三个配对的联合信念相加，就得到该规则的信念 $\pi_t^{+}(h)=\sum_a\Psi_t^{+}(h,a)$；再在该规则下归一化，就得到新的配对权重 $\omega_t^{+}(a\mid h)$。因此，同一次反馈同时改变规则和配对判断，而不是被当作两份独立证据重复使用。

搜索引入新规则时，也需要把配对知识一起迁移。沿用式 \eqref{eq:similarity-transport} 的保留与投影比例：保留分量携带原规则的配对判断，局部投影携带相似旧规则的配对证据，全局投影携带旧工作空间整体的配对信念。这样，更换分类规则不意味着忘掉所有已学的选项关系；新规则的 $\beta$ 则仍从 $\beta_{0,s}$ 开始。

\subsection{为什么部分正确反馈的作用会随学习变化}

在三种配对同样可信时，所选反应的三个其他答案都有可能是同科伙伴。将式 \eqref{eq:c3-feedback-kernel} 对配对取平均，得到规则 $h$ 对两种反馈的概率：
\begin{equation}
\lambda_t(h;r_t=0.5)=\frac{1-p_{t,h}(y_t)}{3},\qquad
\lambda_t(h;r_t=0)=\frac{2[1-p_{t,h}(y_t)]}{3}.
\label{eq:c3-initial-equivalence}
\end{equation}
两者只相差一个与规则无关的倍数，所以在相同先验下，对分类规则产生相同的相对支持。但是，它们对各个配对的支持可以不同。模型由此逐渐积累配对证据，使之后的 $0.5$ 和 0 产生不同作用。

仍以 301 为例，开始时按 4 得到 $0.5$，模型只知道没有完全答对，还不能确定同科伙伴是 1、2 还是 3。后来，如果已有经验使模型相信 2、4 同科，同样的 $0.5$ 就支持“应该考虑按 2”；若得到 0，则支持按 1 或 3。

因此，模型不需要规定“第几个试次以后才使用部分反馈”。变化来自配对信念的逐渐形成，而不是人为切换学习阶段。上面的等价关系只在配对权重均匀时严格成立；相信某个配对，也不等于已经学会如何根据刺激分科。这两种学习结果仍应分别检查。

\subsection{规则精度与搜索怎样调整}
\label{sec:condition3-beta}

部分反馈还需要影响动态规则精度。原有式 \eqref{eq:beta-evidence} 在 $r_t=0.5$ 时更新量恒为零，无法区分该结果是否支持当前规则。第一版保留原来的精度增长和衰减方式，只改变用于判断支持程度的量。

先把当前规则下的三个 $[\Psi_t^{-}(h,a)]^{\gamma_s}$ 归一化，得到经过遗忘、尚未使用当前反馈的配对权重 $\widetilde\omega_t^{-}(a\mid h)$。按这些权重平均反馈概率：
\begin{equation}
\lambda_t(h)=\sum_{a=1}^{3}\widetilde\omega_t^{-}(a\mid h)\ell_t(h,a).
\label{eq:c3-marginal-feedback}
\end{equation}
四个答案等可能时，完全正确、部分正确和零分的机会概率分别为 $1/4$、$1/4$ 和 $1/2$，记为 $b_{\mathrm{chance}}(r_t)$。用规则预测相对于这个机会水平的支持量
\begin{equation}
\kappa_t^{(3)}(h)=
\frac{\lambda_t(h)-b_{\mathrm{chance}}(r_t)}
{\lambda_t(h)+b_{\mathrm{chance}}(r_t)}
\label{eq:c3-beta-support}
\end{equation}
替换式 \eqref{eq:beta-update} 中的 $\kappa_t(h)$。结果比机会水平更符合规则预期时，精度增加；反之，精度降低。这里比较的是实际反馈概率，不能先把它在不同规则之间归一化。更新仍使用反馈前的预测与配对信念，新的 $\beta$ 从下一试次起影响选择。

这一调整不是“得到半分就增加精度”。在配对均匀时，$0.5$ 和 0 得到相同的支持量；在四类等概率时，三种反馈的支持量都为零。该精度更新是需要检验的新假设。

搜索方面，第一版把 0 和 $0.5$ 都视为尚未完全答对，即只在搜索控制公式中用 $z_t=\ind(r_t=1)$ 替换原来的 $r_t$。失败累积、搜索概率和局部／全局范围沿用前文结构。两类反馈仍通过不同的规则后验影响后续局部搜索，原始三值反馈始终保留。

“配对可靠后主动保留科、修改种”暂作为后续比较的机制。先检查上述学习更新是否已能产生合理的科内修订，再判断是否需要额外规定搜索方向。模型不会仅因配对权重集中，就认定分科规则已经掌握。

\subsection{拟合方式与解释范围}

逐试次顺序与基础模型一致：先根据已有经验预测选择，随后观察反应和反馈，再更新规则信念、配对和精度。PF 的每条潜在路径分别保存这些状态，并继续按实际选择的预测概率加权。自主生成时，任务环境按真实答案返回分数：种正确得 1，同科不同种得 $0.5$，不同科得 0；真实配对只供环境评分，不预先告诉模型。

个体参数估计沿用第 \ref{sec:individual-estimation} 节的选择 NLL 目标和搜索流程。配对权重从均匀状态开始、随每次候选运行重新学习，不作为额外的被试参数拟合。四分类距离可以复用，新增计算主要用于配对权重的更新；实际计算开销和所需 PF 预算仍需检查。

结果应分别呈现选择拟合、配对信念以及科与种的学习表现。科正确率、种正确率和平均得分不是同一个指标；对某种配对很确定，也不表示这个配对一定正确。现有固定标签规则库仍需检查是否因反应编码而产生偏向，不能把换编号带来的差异解释为个体策略差异。

本节保留基础模型的工作空间、执行结构和参数种类，主要扩展了配对学习与三值反馈的使用方式。程序已通过实现测试和短序列检查，正式个体拟合与恢复验证尚未开展；具体程序位置、数据核查与检查项目见配套说明 \path{model_0826_condition3_design.md}。


\clearpage
\begin{figure}[p]
\centering
\includegraphics[width=\linewidth]{model_0826_framework.png}
\caption{生成式有限规则搜索模型的逐试次因果结构。作答前，每个粒子依次生成内部知觉、搜索状态、有限工作空间和作答前规则信念，再由动态规则精度和被试层执行结构产生选择概率。PF 对多条不可见路径进行加权平均，得到边际预测。观察实际选择后，PF 先更新粒子权重；随后，选择与反馈共同更新规则证据、有限记忆、所有活跃规则的动态规则精度以及下一试次的搜索状态。当前反馈不能反向改变本试次已经输出的预测，虚线表示跨试次递归。图中为二值反馈基础流程；condition 3 还需第 \ref{sec:condition3-adaptation} 节定义的配对状态和三值反馈更新，图中未单独展开这些扩展部分。}
\label{fig:model-framework}
\end{figure}

\end{document}
